% TOP_PROBLEM: {"id": 1200012, "title": "Some Questions — Question 12", "category_id": 3, "category": {"id": 3, "name": "graph_theory", "display_name": "Graph Theory", "description": "Problems involving graphs, networks, and their properties.", "slug": "graph-theory", "order_index": 3, "created_at": "2026-07-31T15:26:25.671Z"}, "status": "partially_solved", "problem_number": "AMR-011-0012", "set_id": 13}
% FIRST_POSTED: 2026-10-02
% ENTRY_KIND: statement_only
% UnsolvedMath ID 1200012; source catalogue code AMR-011-0012
% !TeX program = lualatex
% Definition and sources only; this file is not an LLM solution attempt.
\documentclass[11pt,a4paper]{article}
\usepackage[margin=25mm]{geometry}
\usepackage{fontspec}
\setmainfont{lmroman10-regular.otf}[BoldFont=lmroman10-bold.otf,
 ItalicFont=lmroman10-italic.otf,BoldItalicFont=lmroman10-bolditalic.otf]
\usepackage{amsmath,amssymb,mathtools}
\usepackage{unicode-math}
\setmathfont{latinmodern-math.otf}
\AtBeginDocument{\let\setminus\smallsetminus}
\usepackage{xurl,hyperref}
\hypersetup{colorlinks=true,urlcolor=blue}
\setcounter{secnumdepth}{0}
\setlength{\parindent}{0pt}
\setlength{\parskip}{5pt}
\setlength{\emergencystretch}{3em}
\begin{document}
\section*{Some Questions — Question 12}\label{1200012}
{\small UnsolvedMath ID 1200012 · AMR-011-0012 · Graph Theory · AMR Open Problem Lists}
\subsection{Definitions and mathematical statement}
Fix a prime $p$. Consider finite undirected graphs $G_j$ of uniformly bounded degree, with an integer label on each edge and, if needed, an integer diagonal label at each vertex. Let $A_j$ be the corresponding symmetric integer adjacency matrix: an off-diagonal entry is the edge label, or zero for a missing edge, and a diagonal entry is the vertex label.

Root $G_j$ at a uniformly selected vertex. Local convergence means that for every radius $r$, the distribution of the rooted, labeled radius-$r$ neighborhood converges to a probability distribution on such rooted neighborhoods, consistently as $r$ varies. Labels are retained in these distributions, not discarded.

Must the normalized ranks after reducing all entries modulo $p$ converge?
\[
 \lim_{j\to\infty}\frac{\operatorname{rank}_{\mathbb F_p}(A_j)}{|V(G_j)|}
                    \quad\hbox{exists?}
\]
Equivalently, ask whether the normalized dimensions of the kernels over the finite field $\mathbb F_p$ converge, since rank plus nullity equals the matrix size.

The matrices are sparse because degrees are bounded, but the dependence among their rows is a global algebraic feature. Real eigenvalue convergence therefore does not by itself prove the asserted finite-field rank convergence. The question includes the ordinary adjacency matrices with every edge label equal to one. The prime is fixed throughout a sequence; varying it with the graph would be a different problem.
\subsection{Short English statement}
Does local convergence of bounded-degree integer-labeled graphs force convergence of adjacency-matrix rank over a fixed finite field?
\subsection{Sources}
\begin{itemize}
\item[S1] ULAM.AI and the UnsolvedMath Contributors, supplied problems.json, record 1200012 (AMR-011-0012), AMR Open Problem Lists. Catalogue identity and source transcription; CC BY 4.0.
\url{https://huggingface.co/datasets/ulamai/UnsolvedMath}.
\item[S2] Abert - Some questions (pdf) (2010), Question 12. Original source indexed by AMR-011-0012.
\url{https://www.renyi.hu/~abert/questions.pdf}.
\end{itemize}
\end{document}
